%% toc: P3 25 25.1
\section{Proofs for Vaught's results on theories categorical in an uncountable power}
\label{sec:annex-p03-w03-p03-007-vaught-categorical}

Throughout, \(T\) is the standing theory of
\cref{def:p03-w03-p03-007-setting}. Two facts are used repeatedly.

\emph{Complete types of tuples.} For a tuple \(\bar a\) of length \(n\) in a
model \(\mathcal C\models T\), let \(\operatorname{tp}^{\mathcal C}(\bar a)\)
be the set of \(\varphi\in F_n(T)\) with \(\mathcal C\models\varphi[\bar a]\).
It is a dual prime ideal by \cref{thm:p03-vaught-algebra-types}. If
\(P\in\mathcal P_n(T)\) and \(\bar a\in P(\mathcal C)\), then
\(P=\operatorname{tp}^{\mathcal C}(\bar a)\): the inclusion
\(P\subseteq\operatorname{tp}^{\mathcal C}(\bar a)\) is the definition of
\(P(\mathcal C)\), and if \(\varphi\notin P\) then \(\neg\varphi\in P\), so
\(\varphi\notin\operatorname{tp}^{\mathcal C}(\bar a)\). If
\(\mathcal C\preccurlyeq\mathcal C'\), a tuple from \(\mathcal C\) has the
same complete type in both, so \(P(\mathcal C)=P(\mathcal C')\cap(\dom\mathcal C)^n\).
An isomorphism \(F\colon\mathcal C\to\mathcal C'\) preserves the truth of
every formula, so \(F[P(\mathcal C)]=P(\mathcal C')\) and
\(F[\theta(\mathcal C)]=\theta(\mathcal C')\) for \(\theta\in F_1(T)\).

\emph{Atoms in complete types.} If \(P\in\mathcal P_n(T)\) contains an atom
\(\alpha\), then \(P\) is principal. Indeed, if \(\varphi\in P\), then
\(\alpha\land\varphi\in P\) is consistent, so \(\vdash_T\alpha\to\varphi\) by
the definition of an atom; conversely \(\vdash_T\alpha\to\varphi\) puts
\(\varphi\) in \(P\) by upward closure. Thus
\(P=\{\varphi:\vdash_T\alpha\to\varphi\}\).

\begin{annexitem}{annex:p03-w03-p03-007-categorical-complete}{thm:p03-w03-p03-007-categorical-complete}{Two models of one power}
\begin{proof}
Suppose that \(T_0\) is not complete, and let \(\sigma\) be a sentence such
that some model \(\mathcal M_1\) of \(T_0\) satisfies \(\neg\sigma\) and some
model \(\mathcal M_2\) of \(T_0\) satisfies \(\sigma\). Put
\(T_1=T_0\cup\{\neg\sigma\}\) and \(T_2=T_0\cup\{\sigma\}\). Each \(T_i\) has
the model \(\mathcal M_i\), which is infinite because \(T_0\) has no finite
model. The language is at most countable and \(\kappa\) is infinite, so
\(\kappa\geq\max(\abs L,\aleph_0)\). By \cref{thm:upward-ls}, each \(T_i\) has
a model \(\mathcal N_i\) of cardinality exactly \(\kappa\). Both
\(\mathcal N_1\) and \(\mathcal N_2\) are models of \(T_0\) of power
\(\kappa\), so they are isomorphic by \(\kappa\)-categoricity. An isomorphism
preserves the truth of every sentence, yet \(\mathcal N_1\models\neg\sigma\)
and \(\mathcal N_2\models\sigma\). This contradiction shows that \(T_0\) is
complete.
\end{proof}
\end{annexitem}

\begin{annexitem}{annex:p03-w03-p03-007-categorical-small}{thm:p03-w03-p03-007-categorical-small}{Both cases of the categorical type count}
\begin{proof}
Suppose first that \(\kappa=\aleph_0\). Then all countably infinite models
of \(T\) are isomorphic, which is clause~1 of
\cref{thm:p03-vaught-finite-formulas-categoricity}. Its clause~3 then says
that every \(\mathcal P_n(T)\) is finite, hence countable.

Suppose next that \(\kappa\) is uncountable. The standing \(T\) is a
complete theory in a countable language with an infinite model. By
\cref{thm:p03-morley-thin-models}, a countable theory with an infinite model
that is categorical in some power \(\kappa>\aleph_0\) is \(\omega\)-stable.
By \cref{thm:p03-vaught-stability-boundary}, an \(\omega\)-stable complete
theory in a countable language is small, that is, every \(\mathcal P_n(T)\)
is countable. This covers every infinite \(\kappa\).
\end{proof}
\end{annexitem}

\begin{annexitem}{annex:p03-w03-p03-007-five-two}{thm:p03-w03-p03-007-five-two}{From countable type sets to the two models}
\begin{proof}
Under the hypothesis of \cref{thm:p03-w03-p03-007-ehrenfeucht}, that cited
theorem gives that every \(\mathcal P_n(T)\) is countable. If instead \(T\) is
\(\kappa\)-categorical for some infinite \(\kappa\),
\cref{thm:p03-w03-p03-007-categorical-small} gives the same conclusion with
a proof. From here on only this conclusion is used.

So \(T\) satisfies clause~1 of
\cref{thm:p03-vaught-smallness-equivalences}, and clause~2 of that theorem
gives a countably infinite saturated model \(\mathcal S\). By the implication
from clause~1 to clause~2 of
\cref{thm:p03-vaught-saturation-characterization}, \(\mathcal S\) is
\(\aleph_0\)-universal and homogeneous. By
\cref{thm:p03-vaught-universal-prime}, since \(T\) has an
\(\aleph_0\)-universal model, it has a prime model.

For uniqueness, let \(\mathcal S'\) be another \(\aleph_0\)-universal,
homogeneous model. It is countably infinite by the definition of
\(\aleph_0\)-universality, and by the implication from clause~2 to clause~1
of \cref{thm:p03-vaught-saturation-characterization} it is saturated. Then
\(\mathcal S'\cong\mathcal S\) by \cref{thm:p03-vaught-saturated-uniqueness}.
Prime models are countably infinite and any two are isomorphic, by
\cref{thm:p03-vaught-prime-atomic}.
\end{proof}
\end{annexitem}

\begin{annexitem}{annex:p03-w03-p03-007-chains}{thm:p03-w03-p03-007-chains}{All seven chain clauses}
\begin{proof}
First, \(\preccurlyeq\) is transitive. If
\(\mathcal C_0\preccurlyeq\mathcal C_1\preccurlyeq\mathcal C_2\), then
\(\mathcal C_0\) is a substructure of \(\mathcal C_2\), and for a formula
\(\varphi\) and a tuple \(\bar a\) from \(\mathcal C_0\),
\(\mathcal C_0\models\varphi[\bar a]\) iff
\(\mathcal C_1\models\varphi[\bar a]\) iff
\(\mathcal C_2\models\varphi[\bar a]\).

Clause~1. We prove by induction on \(\zeta<\gamma\) that
\(\mathcal A_\xi\preccurlyeq\mathcal A_\zeta\) for all \(\xi<\zeta\).
There is nothing to prove for \(\zeta=0\). If \(\zeta=\zeta_0+1\) and
\(\xi<\zeta\), then \(\xi\le\zeta_0\). Either \(\xi=\zeta_0\), or
\(\mathcal A_\xi\preccurlyeq\mathcal A_{\zeta_0}\) by the induction
hypothesis. In both cases
\(\mathcal A_\xi\preccurlyeq\mathcal A_{\zeta_0}\preccurlyeq\mathcal A_\zeta\),
and transitivity applies. If \(\zeta\) is a limit, the induction hypothesis
says that \(\langle\mathcal A_\xi:\xi<\zeta\rangle\) is an elementary chain
in the sense of \cref{thm:elementary-chain}. Since
\(\mathcal A_\zeta\) is its union, that theorem gives
\(\mathcal A_\xi\preccurlyeq\mathcal A_\zeta\) for every \(\xi<\zeta\).

Clause~2. The first assertion is \cref{thm:elementary-chain}. Since
\(\delta>0\), \(\mathcal A_0\preccurlyeq\mathcal A\); every sentence of \(T\)
holds in \(\mathcal A_0\), hence in \(\mathcal A\).

For the remaining clauses, note that a finite tuple \(\bar a\) from
\(\mathcal A\) lies in a single \(\mathcal A_\xi\). Each entry lies in some
\(\mathcal A_{\xi_i}\), and since the universes increase along the chain,
all lie in \(\mathcal A_\xi\) for the largest \(\xi_i\). The empty tuple lies
in \(\mathcal A_0\).

Clause~3. Let \(\bar a\) be a tuple of length \(n\) from \(\mathcal A\), and
choose \(\xi\) with \(\bar a\) in \(\mathcal A_\xi\). Since
\(\mathcal A_\xi\) is atomic, \(\mathcal A_\xi\models\alpha[\bar a]\) for an
atom \(\alpha\) of \(F_n(T)\). By clause~2,
\(\mathcal A\models\alpha[\bar a]\). Thus \(\mathcal A\) is atomic.

Clause~4. Weak saturation: for \(P\in\mathcal P_n(T)\), the saturated
\(\mathcal A_0\) has some \(\bar a\in P(\mathcal A_0)\), and
\(\bar a\in P(\mathcal A)\) because
\(P(\mathcal A_0)=P(\mathcal A)\cap(\dom\mathcal A_0)^n\). Extension: let
\(P\in\mathcal P_n(T)\), \(Q\in\mathcal P_{n+1}(T)\), \(P\subseteq Q\), and
\(\bar a\in P(\mathcal A)\). Choose \(\xi\) with \(\bar a\) in
\(\mathcal A_\xi\). Then \(\bar a\in P(\mathcal A_\xi)\). Since
\(\mathcal A_\xi\) is saturated, some \(b\in\dom\mathcal A_\xi\) has
\((\bar a,b)\in Q(\mathcal A_\xi)\subseteq Q(\mathcal A)\).

Clause~5. A countable union of countable sets is countable under the ambient
Axiom of Choice, so \(\dom\mathcal A\) is countable. It is infinite because
\(\mathcal A\models T\).

Clause~6. If \(x\in\theta(\mathcal A)\), choose \(\xi\) with
\(x\in\dom\mathcal A_\xi\); by clause~2, \(\mathcal A_\xi\models\theta[x]\).
Conversely, \(\mathcal A_\xi\models\theta[x]\) gives
\(\mathcal A\models\theta[x]\), again by clause~2.

Clause~7. The union of \(\aleph_1\) countable sets has power at most
\(\aleph_1\cdot\aleph_0=\aleph_1\). For each \(\xi<\omega_1\) choose
\(x_\xi\in\dom\mathcal A_{\xi+1}\setminus\dom\mathcal A_\xi\). If
\(\xi<\zeta\), then \(x_\xi\in\dom\mathcal A_{\xi+1}\subseteq\dom\mathcal A_\zeta\),
while \(x_\zeta\notin\dom\mathcal A_\zeta\); so \(x_\xi\neq x_\zeta\). The map
\(\xi\mapsto x_\xi\) is an injection of \(\omega_1\) into \(\dom\mathcal A\),
which therefore has power exactly \(\aleph_1\).
\end{proof}
\end{annexitem}

\begin{annexitem}{annex:p03-w03-p03-007-copy}{thm:p03-w03-p03-007-copy}{Fresh elements and the transported structure}
\begin{proof}
Let \(X=\dom\mathcal C\setminus\dom\mathcal D\) and write
\(Z=\dom\mathcal D'\). For \(x\in X\) let \(h(x)\) be the ordered pair
\((Z,x)=\{\{Z\},\{Z,x\}\}\). Then \(h\) is injective. No \(h(x)\) lies in
\(Z\): otherwise \(Z\in\{Z\}\in h(x)\in Z\), a membership cycle forbidden by
the axiom of foundation. Put \(F=g\cup h\). It maps \(\dom\mathcal C\)
bijectively onto \(Z\cup h[X]\), because \(g\) and \(h\) are bijections onto
disjoint sets.

Let \(\mathcal C'\) be the \(L\)-structure on \(Z\cup h[X]\) with
\(R^{\mathcal C'}=\{F(\bar c):\bar c\in R^{\mathcal C}\}\) for each relation
symbol \(R\) and \(c^{\mathcal C'}=F(c^{\mathcal C})\) for each individual
constant. By construction \(F\colon\mathcal C\to\mathcal C'\) is an
isomorphism extending \(g\), and
\(\dom\mathcal C'\setminus Z=F[X]\).

\(\mathcal D'\) is a substructure of \(\mathcal C'\). Each constant
\(c^{\mathcal C}=c^{\mathcal D}\) lies in \(\dom\mathcal D\), so
\(c^{\mathcal C'}=g(c^{\mathcal D})=c^{\mathcal D'}\). For a tuple \(\bar z\)
from \(Z\), write \(\bar z=g(\bar d)\). Then
\(\bar z\in R^{\mathcal C'}\) iff \(\bar d\in R^{\mathcal C}\) iff
\(\bar d\in R^{\mathcal D}\) iff \(\bar z\in R^{\mathcal D'}\).
For elementarity, with \(\bar z=g(\bar d)\) as before and any formula
\(\varphi\),
\[
 \mathcal D'\models\varphi[\bar z]
 \iff\mathcal D\models\varphi[\bar d]
 \iff\mathcal C\models\varphi[\bar d]
 \iff\mathcal C'\models\varphi[F(\bar d)]
 =\varphi[\bar z].
\]
The outer steps use the isomorphisms \(g\) and \(F\), and the middle step
uses \(\mathcal D\preccurlyeq\mathcal C\). The equation
\(\dom\mathcal C'\setminus Z=F[X]\) shows that \(X\) and its image are empty
together. The last assertion is the preservation of \(\theta\)-sets by the
isomorphism \(F\).
\end{proof}
\end{annexitem}

\begin{annexitem}{annex:p03-w03-p03-007-minimal-prime}{thm:p03-w03-p03-007-minimal-prime}{The chain of prime models and the realized nonprincipal type}
\begin{proof}
Suppose, for a contradiction, that a prime model \(\mathcal A\) of \(T\) has
an elementary submodel \(\mathcal A'\) with
\(\dom\mathcal A'\neq\dom\mathcal A\).

\emph{Step 1: \(\mathcal A'\) is prime.} By
\cref{thm:p03-vaught-prime-atomic}, \(\mathcal A\) is countably infinite and
atomic. \(\mathcal A'\) is a model of \(T\), being elementarily equivalent to
\(\mathcal A\), so it is infinite; it is countable as a subset of
\(\dom\mathcal A\). A tuple \(\bar a\) from \(\mathcal A'\) satisfies an atom
\(\alpha\) in \(\mathcal A\), hence in \(\mathcal A'\) by elementarity. So
\(\mathcal A'\) is countably infinite and atomic, and prime by
\cref{thm:p03-vaught-prime-atomic}.

\emph{Step 2: every prime model has a proper prime elementary extension.}
Let \(\mathcal E\) be prime. By \cref{thm:p03-vaught-prime-atomic} there is
an isomorphism \(g\colon\mathcal A'\to\mathcal E\). Apply
\cref{thm:p03-w03-p03-007-copy} to \(\mathcal A'\preccurlyeq\mathcal A\) and
\(g\). It gives \(\mathcal E^+\cong\mathcal A\) with
\(\mathcal E\preccurlyeq\mathcal E^+\) and
\(\dom\mathcal E^+\neq\dom\mathcal E\). As an isomorphic copy of a countably
infinite atomic model, \(\mathcal E^+\) is countably infinite and atomic,
hence prime.

\emph{Step 3: the chain.} By recursion on \(\xi\le\omega_1\), using the
ambient Axiom of Choice to choose the extensions, let
\(\mathcal A_0=\mathcal A\), let
\(\mathcal A_{\xi+1}=(\mathcal A_\xi)^+\) as in Step~2 for \(\xi<\omega_1\),
and let \(\mathcal A_\eta=\bigcup_{\xi<\eta}\mathcal A_\xi\) for limit
\(\eta\le\omega_1\). We show by induction on \(\xi<\omega_1\) that
\(\mathcal A_\xi\) is prime; then Step~2 applies at every successor.
The case \(\xi=0\) is the hypothesis, and successors are handled by Step~2.
At a limit \(\eta<\omega_1\), the sequence below \(\eta\) is an elementary
chain by clause~1 of \cref{thm:p03-w03-p03-007-chains}. Its union is atomic
by clause~3 and countably infinite by clause~5, hence prime by
\cref{thm:p03-vaught-prime-atomic}.

Let \(\mathcal B=\mathcal A_{\omega_1}\). By clauses~1, 2, 3 and~7 of
\cref{thm:p03-w03-p03-007-chains}, \(\mathcal B\) is a model of \(T\) of
power \(\aleph_1\), and it is atomic.

\emph{Step 4: a nonprincipal type.} \(T\) is not \(\aleph_0\)-categorical,
so clause~1 of \cref{thm:p03-vaught-finite-formulas-categoricity} fails, and
with it clause~4: some \(P\in\mathcal P_n(T)\) is nonprincipal. The source
reaches this through its 1.1 and 3.6
\cite[p.~316]{Vaught1961}, which that theorem packages.

\emph{Step 5: a model of power \(\aleph_1\) realizing \(P\).} With new
constants \(c_0,\ldots,c_{n-1}\), the theory \(\overline P\) of
\cref{def:p03-vaught-prime-ideals} is satisfiable by
\cref{thm:p03-vaught-algebra-types}. Its models are expansions of models of
\(T\), so they are infinite, and its language is countable. By
\cref{thm:upward-ls} it has a model \((\mathcal C,c_0,\ldots,c_{n-1})\) of
power \(\aleph_1\). Then \(\mathcal C\models T\) has power \(\aleph_1\) and
\((c_0,\ldots,c_{n-1})\in P(\mathcal C)\).

\emph{Step 6: \(\mathcal B\) omits \(P\).} If \(\bar b\in P(\mathcal B)\),
then \(P=\operatorname{tp}^{\mathcal B}(\bar b)\). Since \(\mathcal B\) is
atomic, \(\bar b\) satisfies an atom \(\alpha\), so \(\alpha\in P\), and \(P\)
would be principal by the fact recalled at the head of this annex. Hence
\(P(\mathcal B)=\varnothing\).

\emph{Step 7.} An isomorphism would carry \(P(\mathcal C)\neq\varnothing\)
onto \(P(\mathcal B)=\varnothing\), so \(\mathcal B\not\cong\mathcal C\).
Both are models of \(T\) of power \(\aleph_1\), which contradicts
\(\aleph_1\)-categoricity. Hence \(\mathcal A\) has no proper elementary
submodel.
\end{proof}
\end{annexitem}

\begin{annexitem}{annex:p03-w03-p03-007-auxiliary}{thm:p03-w03-p03-007-auxiliary}{Relativization, the sentences (1) and (2), and the two interpretations}
\begin{proof}
Clause~1. We prove, by induction on \(L\)-formulas \(\varphi\), that for
every assignment \(s\) of elements of \(W\) to the variables,
\(\mathcal C_W\models\varphi[s]\) iff \((\mathcal C,W)\models\varphi^V[s]\).
The clause is the special case of assignments sending \(v_i\) to \(a_i\) for
\(i<n\), since the truth of \(\varphi\in F_n(T)\) depends only on those
variables. An atomic formula is built from a relation symbol or identity and
from variables and constants. Its relativization is itself, and all its
arguments denote elements of \(W\), because the constants of \(\mathcal C\)
lie in the substructure. The relations of \(\mathcal C_W\) are the
restrictions to \(W\) of those of \(\mathcal C\), so the two truth values
agree. The connective cases are immediate. For
\(\exists v_j\psi\), \((\mathcal C,W)\models\exists v_j(Vv_j\land\psi^V)[s]\)
iff some \(w\in W\) has \((\mathcal C,W)\models\psi^V[s(v_j/w)]\). The
modified assignment still takes values in \(W\), so by the induction
hypothesis this holds iff some \(w\in W\) has
\(\mathcal C_W\models\psi[s(v_j/w)]\), that is,
\(\mathcal C_W\models\exists v_j\psi[s]\). The universal case is dual.

Clause~2. The sentences of \(T\) hold in \(\mathcal C\). In (1),
\(\dom\mathcal D\) is nonempty, it differs from \(\dom\mathcal C\), and
\(\theta(\mathcal C)\subseteq\dom\mathcal D\) by hypothesis. For (2), let
\(\bar a\) be an \(n\)-tuple from \(\dom\mathcal D\) and
\(\varphi\in F_n(T)\). Then \(\mathcal C\models\varphi[\bar a]\) iff
\(\mathcal D\models\varphi[\bar a]\) by elementarity, iff
\((\mathcal C,\dom\mathcal D)\models\varphi^V[\bar a]\) by clause~1 with
\(W=\dom\mathcal D\).

Clause~3. Let \(W=V^{\mathfrak B}\). It is nonempty by the first sentence of
(1). For an individual constant \(c\), the sentence
\(\sigma_c=\exists v_0(v_0=c)\) is in \(F_0(T)\) and is logically valid, so
\(\mathfrak B\models\sigma_c\). The instance of (2) with \(n=0\) gives
\(\mathfrak B\models\sigma_c^V\), that is, \(\exists v_0(Vv_0\land v_0=c)\);
so \(c^{\mathcal B}\in W\). Since \(L\) has no function symbols, \(W\) is
the universe of the substructure \(\mathcal B_V\), and clause~1 applies. For
an \(n\)-tuple \(\bar a\) from \(W\) and \(\varphi\in F_n(T)\),
\[
 \mathcal B\models\varphi[\bar a]
 \iff\mathfrak B\models\varphi^V[\bar a]
 \iff\mathcal B_V\models\varphi[\bar a],
\]
by the instance of (2) for \(\varphi\) and by clause~1. Every \(L\)-formula
lies in some \(F_m(T)\), and its truth under an assignment into \(W\) is
truth of such a tuple instance. Thus \(\mathcal B_V\preccurlyeq\mathcal B\).
\(\mathcal B\models T\) because \(T\subseteq T'\), and \(\mathcal B_V\models T\)
by elementarity. The second sentence of (1) gives an element outside \(W\),
so the submodel is proper. By the third sentence of (1),
\(\theta(\mathcal B)\subseteq W\). For \(x\in W\),
\(\mathcal B\models\theta[x]\) iff \(\mathcal B_V\models\theta[x]\). Hence
\(\theta(\mathcal B)=\theta(\mathcal B_V)\).

Clause~4. Both domain formulas are nonempty in every model of \(T'\): for
\(v_0=v_0\) this is logical, and for \(Vv_0\) it is (1). Each constant is
interpreted by a target constant, and that constant lies in the domain,
trivially for \(I_=\) and by clause~3 for \(I_V\). It remains to check that
\(T'\models\sigma^I\) for each \(\sigma\in T\). For \(I_V\), \(\sigma^{I_V}\)
is \(\sigma^V\), and the instance \(\sigma\leftrightarrow\sigma^V\) of (2),
together with \(\sigma\in T'\), gives \(T'\models\sigma^V\). For \(I_=\),
\(\sigma^{I_=}\) differs from \(\sigma\) only by conjoining or premising the
valid \(v_j=v_j\) to each quantifier, so it is logically equivalent to
\(\sigma\in T'\). The denotations are as stated: \(I_=(\mathfrak B)\) has
universe \(\dom\mathcal B\) and the relations of \(\mathcal B\), and
\(I_V(\mathfrak B)\) has universe \(V^{\mathfrak B}\) and the restricted
relations. Finally, \(L_V\) adds one symbol to the countable \(L\).
\end{proof}
\end{annexitem}

\begin{annexitem}{annex:p03-w03-p03-007-definable-sets}{thm:p03-w03-p03-007-definable-sets}{Vaught's two-model argument}
\begin{proof}
Suppose, for a contradiction, that \(\mathcal A^1\models T\) and that
\(U=\theta(\mathcal A^1)\) is infinite with \(\abs U<\abs{\dom\mathcal A^1}\).

\emph{Step 1: \(T'\) is consistent.} By \cref{thm:downward-ls}, there is
\(\mathcal A^2\preccurlyeq\mathcal A^1\) with \(U\subseteq\dom\mathcal A^2\)
and \(\abs{\dom\mathcal A^2}\le\max(\abs U,\abs L,\aleph_0)=\abs U\), the
last equality because \(U\) is infinite and \(L\) countable. So
\(\abs{\dom\mathcal A^2}=\abs U<\abs{\dom\mathcal A^1}\), and
\(\mathcal A^2\) is a proper elementary submodel containing
\(\theta(\mathcal A^1)\). By clause~2 of
\cref{thm:p03-w03-p03-007-auxiliary},
\((\mathcal A^1,\dom\mathcal A^2)\models T'\).

\emph{Step 2: \(\theta\) is infinite in every model of \(T\).} For
\(k\ge1\), let \(\sigma_k\) be the sentence
\(\exists v_1\cdots\exists v_k
 (\bigwedge_{1\le i<j\le k}v_i\neq v_j\land\bigwedge_{i=1}^k\theta(v_i))\),
where \(\theta(v_i)\) is the proper substitution of \(v_i\) for \(v_0\).
Each \(\sigma_k\) is true in \(\mathcal A^1\). Since \(T\) is complete, every
\(\sigma_k\) belongs to \(T\), so \(\theta(\mathcal C)\) is infinite for every
\(\mathcal C\models T\).

\emph{Step 3: \(T\) is small.} By
\cref{thm:p03-w03-p03-007-categorical-small} with \(\kappa=\aleph_1\), every
\(\mathcal P_n(T)\) is countable. The source cites
\cref{thm:p03-w03-p03-007-ehrenfeucht} here.

\emph{Step 4: two saturated interpretations.} \(T'\) is consistent and
countable, and \(I_=\), \(I_V\) are relative interpretations of \(T\) in
\(T'\), by clause~4 of \cref{thm:p03-w03-p03-007-auxiliary}. By
\cref{thm:p03-vaught-saturated-existence}, with the family
\(\{I_=,I_V\}\), there is a countably infinite \(\mathfrak B\models T'\) such
that \(\mathcal B=I_=(\mathfrak B)\) and \(\mathcal B_V=I_V(\mathfrak B)\)
are saturated models of \(T\). By clause~3 of
\cref{thm:p03-w03-p03-007-auxiliary}, \(\mathcal B_V\) is a proper elementary
submodel of \(\mathcal B\) with \(\theta(\mathcal B)=\theta(\mathcal B_V)\).
Both are countably infinite: \(\mathcal B\) by construction, and
\(\mathcal B_V\) as an infinite model of \(T\) inside it. This is the source's
display (3).

\emph{Step 5: every countable saturated model has such an extension.} Let
\(\mathcal S\) be a countably infinite saturated model of \(T\). By
\cref{thm:p03-vaught-saturated-uniqueness} there is an isomorphism
\(g\colon\mathcal B_V\to\mathcal S\). Apply \cref{thm:p03-w03-p03-007-copy}
to \(\mathcal B_V\preccurlyeq\mathcal B\) and \(g\). It gives
\(\mathcal S^*\) and an isomorphism \(F\colon\mathcal B\to\mathcal S^*\)
extending \(g\), with \(\mathcal S\) a proper elementary submodel of
\(\mathcal S^*\). Saturation is preserved by isomorphisms, since they carry
each \(P(\cdot)\) onto the corresponding set; so \(\mathcal S^*\) is a
countably infinite saturated model. Moreover
\[
 \theta(\mathcal S^*)=F[\theta(\mathcal B)]=F[\theta(\mathcal B_V)]
 =g[\theta(\mathcal B_V)]=\theta(\mathcal S).
\]
Write \(\mathcal S^*\) for one such extension, chosen by the ambient Axiom of
Choice.

\emph{Step 6: the chain.} By clause~2 of
\cref{thm:p03-vaught-smallness-equivalences} and Step~3, there is a countably
infinite saturated \(\mathcal A_0\models T\). Define
\(\mathcal A_{\xi+1}=(\mathcal A_\xi)^*\) for \(\xi<\omega_1\), and
\(\mathcal A_\eta=\bigcup_{\xi<\eta}\mathcal A_\xi\) for limit
\(\eta\le\omega_1\). By induction on \(\xi<\omega_1\), every
\(\mathcal A_\xi\) is countably infinite and saturated, and
\(\theta(\mathcal A_\xi)=\theta(\mathcal A_0)\). Successor stages are
Step~5. At a limit \(\eta<\omega_1\), the earlier stages form an elementary
chain by clause~1 of \cref{thm:p03-w03-p03-007-chains}. Clauses~4 and~5
make the union countably infinite and saturated, and clause~6 gives
\(\theta(\mathcal A_\eta)=\bigcup_{\xi<\eta}\theta(\mathcal A_\xi)=\theta(\mathcal A_0)\).
This is the source's remark that the union of an elementarily increasing
sequence of saturated systems is saturated \cite[p.~318]{Vaught1961}.

Let \(\mathcal A_{\omega_1}\) be the union of all the stages. By clauses~1,
2, 6 and~7 of \cref{thm:p03-w03-p03-007-chains}, it is a model of \(T\) of
power \(\aleph_1\) with
\(\theta(\mathcal A_{\omega_1})=\theta(\mathcal A_0)\). This set is countable
because \(\mathcal A_0\) is, and infinite by Step~2, so it has power
\(\aleph_0\).

\emph{Step 7: a model whose \(\theta\)-set has power \(\aleph_1\).} Add to
\(L\) new constants \(c_\xi\) for \(\xi<\omega_1\), and let \(T_1\) consist of
\(T\), the sentences \(c_\xi\neq c_\eta\) for \(\xi\neq\eta\), and the
sentences \(\theta(c_\xi)\). A finite subset mentions finitely many
constants. Interpreting them by distinct elements of the infinite set
\(\theta(\mathcal A_0)\) makes it true in an expansion of \(\mathcal A_0\). By
\cref{thm:fol-compactness}, with the Axiom of Choice for this uncountable
language, \(T_1\) has a model \(\mathcal N\). Let \(A\) be the set of
interpretations of the \(c_\xi\), of power \(\aleph_1\). By
\cref{thm:downward-ls} there is \(\mathcal N'\preccurlyeq\mathcal N\) with
\(A\subseteq\dom\mathcal N'\) and \(\abs{\dom\mathcal N'}\le\aleph_1\). So
\(\mathcal N'\) has power \(\aleph_1\). Its \(L\)-reduct \(\mathcal C\) is a
model of \(T\) of power \(\aleph_1\). Each \(c_\xi\) satisfies \(\theta\) in
\(\mathcal N'\), by elementarity, and hence in \(\mathcal C\). Thus
\(A\subseteq\theta(\mathcal C)\subseteq\dom\mathcal C\), and
\(\theta(\mathcal C)\) has power \(\aleph_1\).

\emph{Step 8.} An isomorphism
\(\mathcal A_{\omega_1}\to\mathcal C\) would map
\(\theta(\mathcal A_{\omega_1})\) bijectively onto \(\theta(\mathcal C)\).
These sets have powers \(\aleph_0\neq\aleph_1\), so there is no such
isomorphism. This contradicts \(\aleph_1\)-categoricity. Hence, for every
\(\mathcal A\models T\), \(\theta(\mathcal A)\) is finite, or it is infinite
and not of smaller power than \(\dom\mathcal A\). In the second case it has
the same power, since \(\theta(\mathcal A)\subseteq\dom\mathcal A\) and
cardinals are comparable under the Axiom of Choice.
\end{proof}
\end{annexitem}

\begin{annexitem}{annex:p03-w03-p03-007-morley-transfer}{thm:p03-w03-p03-007-morley-transfer}{Reduction to \(\aleph_1\) by Morley's theorem}
\begin{proof}
The standing \(T\) is complete, in a countable language, with an infinite
model. By \cref{thm:p03-morley-original-categoricity}, categoricity in the
uncountable power \(\kappa\) implies categoricity in every uncountable power,
in particular in \(\aleph_1\). Then
\cref{thm:p03-w03-p03-007-definable-sets} applies. If moreover \(T\) is not
\(\aleph_0\)-categorical, \cref{thm:p03-w03-p03-007-minimal-prime} applies.
\end{proof}
\end{annexitem}

\begin{annexitem}{annex:p03-w03-p03-007-acf}{thm:p03-w03-p03-007-acf}{The field computations and the two model types}
\begin{proof}
By \cref{thm:p03-vaught-functional-bridge}, the ring language may be
replaced by its relational graph presentation. That replacement preserves
models, cardinalities, types, primeness, atomicity, saturation and
homogeneity, so Vaught's results apply to \(\mathsf{ACF}_0\) as stated below.

Clause~1. By \cref{thm:p03-bl-acf-geometry}, \(\mathsf{ACF}_0\) is complete
and strongly minimal. Its countable models are exactly one for each
transcendence degree \(0,1,2,\ldots,\aleph_0\), so there are infinitely many
non-isomorphic countable models and the theory is not
\(\aleph_0\)-categorical. By \cref{thm:p03-bl-whole-models}, a strongly
minimal theory in a countable language is categorical in every uncountable
power. A model is infinite: if \(K=\{a_1,\ldots,a_m\}\) were a finite field,
the polynomial \((x-a_1)\cdots(x-a_m)+1\) would have no root in \(K\).

Clause~2. \emph{Algebraic elements form a subfield.} Let
\(a_1,\ldots,a_m\in K\) be algebraic over \(\mathbb Q\). Each \(a_i\) is a root
of a monic polynomial over \(\mathbb Q\) of some degree \(d_i\ge1\): divide a
nonzero polynomial vanishing at \(a_i\) by its leading coefficient. Let
\(W\) be the \(\mathbb Q\)-linear span of the monomials
\(a_1^{e_1}\cdots a_m^{e_m}\) with \(e_i<d_i\). Multiplying such a monomial by
\(a_i\) gives either another such monomial or one containing
\(a_i^{d_i}\). The monic equation rewrites \(a_i^{d_i}\) as a
\(\mathbb Q\)-combination of lower powers, so \(W\) is closed under
multiplication by each \(a_i\). Hence \(W\) is a subring containing
\(\mathbb Q[a_1,\ldots,a_m]\), of finite dimension \(N\) over \(\mathbb Q\).
For \(c\in W\), the \(N+1\) elements \(1,c,\ldots,c^N\) of \(W\) are linearly
dependent, so \(c\) is a root of a nonzero polynomial over \(\mathbb Q\) and is
algebraic. Thus sums and products of algebraic elements are algebraic. If
\(a\) is algebraic, so is \(-a\), a root of \(p(-x)\) when \(p(a)=0\). If
also \(a\neq0\), cancel powers of \(x\) from \(p\) so that its constant term
is nonzero. Then \(a^{-1}\) is a root of the reversed polynomial
\(x^{\deg p}p(1/x)\), which is nonzero. So \(\mathbb A_K\) is a subfield
containing \(\mathbb Q\).

\emph{It is algebraically closed.} Let
\(f(x)=x^m+e_{m-1}x^{m-1}+\cdots+e_0\) with \(m\ge1\) and
\(e_i\in\mathbb A_K\); a nonconstant polynomial over the field \(\mathbb A_K\)
becomes monic after division by its leading coefficient. Since \(K\) is
algebraically closed, some \(c\in K\) has \(f(c)=0\). As above, let the
\(e_i\) have monic equations of degrees \(d_i\), and let \(W\) be the span of
the products of monomials \(e_0^{j_0}\cdots e_{m-1}^{j_{m-1}}\) (\(j_i<d_i\))
with \(c^j\) (\(j<m\)). This span is closed under multiplication by each
\(e_i\), as before, and by \(c\), since
\(c^m=-(e_{m-1}c^{m-1}+\cdots+e_0)\). It is finite-dimensional and contains
the powers of \(c\), so \(c\) is algebraic over \(\mathbb Q\) and lies in
\(\mathbb A_K\).

\emph{It is countable.} \(\mathbb Q[x]\) is countable, and each nonzero
polynomial has at most as many roots in the field \(K\) as its degree.

\emph{It is prime and minimal.} \(\mathbb A_K\) is an algebraically closed
field of characteristic \(0\), so it is a model of \(\mathsf{ACF}_0\). By
\cref{thm:p03-bl-acf-geometry}, its algebraic closure of the empty set is the
set of its elements algebraic over \(\mathbb Q\), which is all of
\(\mathbb A_K\). So the empty set spans, and it is a basis: the model has
dimension \(0\). By \cref{thm:p03-bl-whole-models}, \(\mathbb A_K\) is prime
over this basis. The empty map into a model \(N\) of \(\mathsf{ACF}_0\) is
elementary, because \(\mathbb A_K\) and \(N\) satisfy the same sentences of
the complete theory. Hence it extends to an elementary embedding
\(\mathbb A_K\to N\), and \(\mathbb A_K\) is prime. Minimality follows from
\cref{thm:p03-w03-p03-007-minimal-prime}, by clause~1. It also follows
directly: an elementary submodel \(M\) of \(\mathbb A_K\) is a model of
\(\mathsf{ACF}_0\) and contains all elements algebraic over \(\mathbb Q\)
inside it. Each \(a\in\mathbb A_K\) is a root of a nonzero polynomial \(p\)
with integer coefficients, obtained by clearing denominators. If \(p\) has
exactly \(t\) roots in \(\mathbb A_K\), the sentence saying that \(p(x)=0\)
has exactly \(t\) solutions is true in \(\mathbb A_K\), hence in \(M\). The
\(t\) roots in \(M\) are among the \(t\) roots in \(\mathbb A_K\), so they
are all of them, and \(a\in M\). Any two such fields have
dimension \(0\), so they are isomorphic by \cref{thm:p03-bl-whole-models}.
For \(K=\mathbb C\), \(\mathbb A_K\) is by definition the field of complex
numbers algebraic over \(\mathbb Q\).

Clause~3. Let \(K\) be countable with transcendence degree \(\aleph_0\). By
\cref{thm:p03-bl-acf-geometry}, a pregeometry basis is a transcendence basis,
so \(K\) has a basis \(I\) with \(\abs I=\aleph_0\). We prove that \(K\) is
saturated in the sense of \cref{def:p03-vaught-saturation}. Strong
minimality, as stated in \cref{def:p03-bl-strongly-minimal-theory}, holds in
every model, and the uniform bound of \cref{thm:p03-bl-uniform} transfers
between models. For \(\varphi(v_0,\ldots,v_n)\), that lemma gives \(q\) such
that the sentence
\[
 \forall v_0\cdots\forall v_{n-1}\,
 \bigl(\exists^{\le q}v_n\,\varphi\lor\exists^{\le q}v_n\,\neg\varphi\bigr)
\]
holds in one model. Here \(\exists^{\le q}\) is the usual first-order
abbreviation. Since \(\mathsf{ACF}_0\) is complete, the sentence holds in all
models. As models are infinite, \(\varphi(\bar a,x)\) defines a cofinite set
in a model exactly when the formula
\(\gamma_\varphi(\bar v)=\exists^{\le q}v_n\,\neg\varphi\) holds at \(\bar a\).

\emph{Extension.} Let \(P\in\mathcal P_n(T)\), \(Q\in\mathcal P_{n+1}(T)\),
\(P\subseteq Q\), and \(\bar a\in P(K)\), with \(n=0\) allowed. By
\cref{thm:p03-vaught-algebra-types}, \(Q\) is realized by some
\((\bar a',b')\) in a model \(N\). Then
\(\operatorname{tp}^N(\bar a')=Q\cap F_n(T)=P=\operatorname{tp}^K(\bar a)\).

Case~1: \(b'\) is algebraic over \(\bar a'\) in \(N\). Choose
\(\varphi\in F_{n+1}(T)\) with \(N\models\varphi[\bar a',b']\) and
\(\varphi(\bar a',x)\) having exactly \(k\ge1\) solutions in \(N\). Then
\(\varphi\in Q\), and the formula
\(\exists^{=k}v_n\varphi\in F_n(T)\) lies in \(P\). So \(\varphi(\bar a,x)\) has
exactly \(k\) solutions \(b_1,\ldots,b_k\) in \(K\). If no \((\bar a,b_i)\)
belonged to \(Q(K)\), choose \(\chi_i\in Q\) false at \((\bar a,b_i)\), and
put \(\chi=\varphi\land\chi_1\land\cdots\land\chi_k\in Q\). Then
\(\exists v_n\chi\in Q\cap F_n(T)=P\), so some \(b\in K\) satisfies
\(\chi(\bar a,b)\). Such a \(b\) satisfies \(\varphi(\bar a,b)\), so it is
some \(b_i\), at which \(\chi_i\) fails, a contradiction. Hence some
\((\bar a,b_i)\in Q(K)\).

Case~2: \(b'\) is not algebraic over \(\bar a'\). For each \(\varphi\), the
set \(\varphi(\bar a',N)\) is finite or cofinite. If it is finite, \(b'\) is
not in it; if it is cofinite, \(b'\) is not in its finite complement. So
\(\varphi\in Q\) iff \(\varphi(\bar a',x)\) is cofinite in \(N\), iff
\(\gamma_\varphi\in P\). The set \(I\) is independent and infinite, while an
independent set inside the closure of the \(n\) entries of \(\bar a\) has at
most \(n\) elements by \cref{thm:p03-bl-dimension}. So some \(b\in I\) is
not algebraic over \(\bar a\) in \(K\). By the same reasoning in \(K\),
\(K\models\varphi[\bar a,b]\) iff \(\gamma_\varphi\in\operatorname{tp}^K(\bar a)=P\)
iff \(\varphi\in Q\). Thus \((\bar a,b)\in Q(K)\).

\emph{Weak saturation.} Let \(P\in\mathcal P_n(T)\) and
\(P_k=P\cap F_k(T)\) for \(k\le n\). Each \(P_k\) is a dual prime ideal of
\(F_k(T)\): it contains the valid formulas, omits contradictions, and
inherits closure and the decision property from \(P\). Also
\(P_k\subseteq P_{k+1}\). The empty tuple lies in \(P_0(K)\) because
\(K\models\mathsf{ACF}_0\) and \(P_0\) is the set of sentences true in the
models of the complete theory. Applying the extension property \(n\) times
gives a tuple in \(P_n(K)=P(K)\). So \(K\) is saturated, and by
\cref{thm:p03-vaught-saturation-characterization} it is
\(\aleph_0\)-universal and homogeneous.

\emph{Existence and the converse.} By \cref{thm:p03-bl-acf-geometry} there
is a countable model of transcendence degree \(\aleph_0\), and it is unique
up to isomorphism. Let \(K'\) be a countable model that is saturated, or
equivalently, by \cref{thm:p03-vaught-saturation-characterization}, a
countably infinite model that is \(\aleph_0\)-universal and homogeneous. By
\cref{thm:p03-vaught-saturated-uniqueness}, \(K'\) is isomorphic to the
saturated model of transcendence degree \(\aleph_0\). A field isomorphism
fixes the prime field and carries transcendence bases to transcendence
bases, so \(K'\) has transcendence degree \(\aleph_0\). This proves the
three-way equivalence and the uniqueness.
\end{proof}
\end{annexitem}

\begin{annexitem}{annex:p03-w03-p03-007-real-types}{thm:p03-w03-p03-007-real-types}{Rational cuts separate the types of real numbers}
\begin{proof}
\emph{Upper bound.} The language \(\{0,1,+,-,\cdot\}\) is countable, and a
formula is a finite string over a countable alphabet, so the set of
formulas with free variables among \(v_0\) is countable. A complete one-type
is a subset of that set. Hence there are at most \(2^{\aleph_0}\) complete
one-types.

\emph{Numerals and the separating formulas.} For an integer \(m\), let
\(\underline m\) be the term \(0\) when \(m=0\), the sum
\(1+\cdots+1\) of \(m\) summands when \(m>0\), and
\(-(1+\cdots+1)\) with \(\abs m\) summands when \(m<0\). In \(\mathbb R\) it
denotes \(m\). For an integer \(m\) and a positive integer \(k\), put
\[
 \psi_{m,k}(v_0)\;=\;\exists v_1\,
   \bigl(\underline k\cdot v_0+v_1\cdot v_1=\underline m\bigr).
\]
For \(x\in\mathbb R\), \(\mathbb R\models\psi_{m,k}[x]\) if and only if
\(m-kx\) is the square of a real number. Squares of reals are nonnegative,
and every nonnegative real is the square of its nonnegative square root.
Hence \(\mathbb R\models\psi_{m,k}[x]\) if and only if \(m-kx\ge0\), that
is, if and only if \(x\le m/k\).

\emph{Injectivity.} For \(x\in\mathbb R\), the set
\(\operatorname{tp}(x)=\{\varphi(v_0):\mathbb R\models\varphi[x]\}\) is a
complete one-type of \(T_{\mathbb R}\): it is realized in the model
\(\mathbb R\) of \(T_{\mathbb R}\), and it contains \(\varphi\) or
\(\neg\varphi\) for each \(\varphi\). Let \(r<s\) be real. By the
Archimedean property choose a positive integer \(k\) with \(k(s-r)>1\), and
let \(m\) be the least integer with \(m\ge kr\). Then \(m<kr+1<ks\), so
\(r\le m/k<s\). Hence \(\psi_{m,k}\in\operatorname{tp}(r)\) and
\(\neg\psi_{m,k}\in\operatorname{tp}(s)\), and
\(\operatorname{tp}(r)\neq\operatorname{tp}(s)\).

\emph{Count.} The map \(x\mapsto\operatorname{tp}(x)\) is an injection of
\(\mathbb R\), of power \(2^{\aleph_0}\), into the set of complete
one-types. With the upper bound, the Cantor--Schr\"oder--Bernstein theorem
gives exactly \(2^{\aleph_0}\) complete one-types.

\emph{Graph language.} By \cref{thm:p03-vaught-functional-bridge}, formula
translation into the graph language of \(\{+,-,\cdot\}\), with \(0\) and
\(1\) kept as constants, identifies the complete types of \(T_{\mathbb R}\)
with those of its graph-language presentation, and the latter are the
members of \(\mathcal P_1\) by \cref{thm:p03-vaught-algebra-types}. So
\(\mathcal P_1\) has \(2^{\aleph_0}\) members.
\end{proof}
\end{annexitem}

\begin{annexitem}{annex:p03-w03-p03-007-pure-sets}{thm:p03-w03-p03-007-pure-sets}{Categoricity of pure sets and a proper elementary submodel}
\begin{proof}
\(T_\infty\) is complete, and its prime models are exactly the countably
infinite ones, by \cref{thm:p03-vaught-finite-boundary}. In the language
with identity only, a bijection between two sets of the same power is an
isomorphism, so \(T_\infty\) is categorical in every infinite power.

Let \(M=\omega\) and \(M'=\omega\setminus\{0\}\). With no relation symbols or
constants, \(M'\) is a substructure. We use the Tarski--Vaught test,
\cref{thm:tarski-vaught}. Let \(M\models\exists y\,\varphi(y,\bar a)\) with
\(\bar a\) from \(M'\), and let \(b\in M\) be a witness. If \(b\neq0\), then
\(b\in M'\). If \(b=0\), choose \(b'\in M'\) different from every entry of
\(\bar a\). The permutation of \(\omega\) exchanging \(0\) and \(b'\) and
fixing everything else is an automorphism. It fixes \(\bar a\), so
\(M\models\varphi(b',\bar a)\) with \(b'\in M'\). Hence \(M'\preccurlyeq M\),
and \(M'\neq M\). The theory is \(\aleph_1\)-categorical, but it is also
\(\aleph_0\)-categorical, and the conclusion of
\cref{thm:p03-w03-p03-007-minimal-prime} fails for it.
\end{proof}
\end{annexitem}

\begin{annexitem}{annex:p03-w03-p03-007-two-cardinal-boundary}{thm:p03-w03-p03-007-two-cardinal-boundary}{The structure \(\mathcal M\) and the contrapositive}
\begin{proof}
\(T_U\) is the set of sentences true in \(\mathcal M\), so it is complete,
its language is countable, and it has the infinite model \(\mathcal M\). Here
\(\theta(\mathcal M)=\omega\) is infinite and has power
\(\aleph_0<\aleph_1=\abs{\omega_1}\). If \(T_U\) were \(\aleph_1\)-categorical,
\cref{thm:p03-w03-p03-007-definable-sets} would make \(\theta(\mathcal M)\)
finite or of power \(\aleph_1\). Hence \(T_U\) is not
\(\aleph_1\)-categorical, and the conclusion of that theorem fails for a
complete theory that lacks its hypothesis.
\end{proof}
\end{annexitem}
